\chapter{Geltungsbereich und ZF-Axiome}

Dieser Referenzband sammelt grundlegende Definitionen und Hauptresultate
über Mengen ohne zusätzliche Struktur. Seine Grundlage ist
\emph{ZF}: die Zermelo--Fraenkel-Mengenlehre ohne Auswahlaxiom.
Das Auswahlaxiom wird weder vorausgesetzt noch verneint.
Alle aufgenommenen Sätze folgen bereits aus ZF.

Der Leitfaden ergänzt Band~3 und folgt dem beweisfreien Format der
Leitfäden zu Band~28 und Band~37. Er enthält eine ausgewählte Übersicht
bekannter Resultate, keine vollständige Aufzählung aller
mengentheoretischen Sätze. Beweise und Beweisskizzen entfallen.

\section{Abgrenzung in der Bandfolge}

Die Zuordnung wurde mit dem Überblicksband, den Fachbandquellen und
dem Bandverzeichnis abgeglichen. Ein in späteren Bänden lediglich
benutzter Mengenbegriff bleibt hier zulässig; eigenständige Begriffe
und Theorien dieser Bände werden nicht entwickelt.

\begin{center}
\small
\begin{tabularx}{\linewidth}{@{}p{.59\linewidth}X@{}}
\toprule
Ausgelassene Themen & Zuständige Bände\\
\midrule
Logische Kalküle und allgemeine logische Sätze & 1--2\\
Relationen, Funktionen sowie Injektivität, Surjektivität und Bijektivität
  & 4--8; Grundbegriffe bereits in 3\\
Auswahlprinzipien und Wohlordnungssatz & 9 und 16\\
Natürliche Zahlen, Zahlbereiche und Folgen & 10, 17--19 und 21\\
Äquivalenzrelationen, Quotienten und Gleichmächtigkeit & 11\\
Ordnungen, Schranken, Infima und Suprema & 12--16\\
Endlichkeit, Anzahlen und endliche Kombinatorik & 20\\
Graphen, Bäume und Wörter & 22--27\\
Algebraische Strukturen und ihre Verknüpfungen & 28--45\\
Spezielle vereinigungsabgeschlossene Familien und Frankls Vermutung & 46\\
Metrische und topologische Begriffe & 47\\
Ordinale, Kardinale, transitive Mengen, Ränge, Rekursion,
Modelle und Unabhängigkeit & 48\\
\bottomrule
\end{tabularx}
\end{center}

Geordnete Paare und kartesische Produkte werden zwar bereits in Band~3
konstruiert, bleiben hier als Werkzeuge der anschließenden
Relations- und Funktionstheorie ebenfalls ausgenommen.
Partitionen, Auswahlmengen und spezielle abgeschlossene
Mengensysteme sind kein eigener Gegenstand dieses Leitfadens.

\section{Sprache und Leseregeln}

Alle Objektvariablen bezeichnen Mengen; auch die Elemente einer Menge
sind Mengen. Die Grundzeichen sind Gleichheit und Mitgliedschaft
$\in$. Logische Zeichen werden im Sinn der Bände~1 und~2 verwendet.
Ungebundene Mengenvariablen in einer Aussage sind allgemein gemeint.
Es werden keine zusätzlichen Strukturaxiome an Mengen gestellt.
Insbesondere wird keine Endlichkeit der vorgegebenen Mengen
vorausgesetzt. Funktionen und Relationen werden nicht als zusätzliche
Objekte vorausgesetzt.
Davon zu unterscheiden sind ausdrücklich genannte Bedingungen innerhalb
eines Satzes: Eine Implikation $A\subseteq B\Rightarrow\cdots$ gilt
für alle Mengen $A,B$; ihre Folgerung allein braucht nicht immer zu gelten.
Der absolute Durchschnitt $\bigcap\mathcal F$ ist nur für
$\mathcal F\ne\varnothing$ definiert. Der relative Durchschnitt
$\bigcap_U\mathcal F$ schließt dagegen leere Familien ein.
Auch die Eindeutigkeitsbedingung im Ersetzungsschema und die
Existenzbedingung im Sammlungsschema gehören zur jeweiligen Aussage.

$\varphi$ und $\psi$ bezeichnen Formeln der Mengensprache mit
festen Mengenparametern. Bei einem \emph{Schema} steht für jede
solche Formel eine eigene Aussage. Diese Schreibweise führt
keine Mengen von Formeln oder weiteren mathematischen Strukturen ein.
In den Schemata ist die jeweils neu behauptete Menge keine freie
Variable der eingesetzten Formel.

\begin{definition}[Teilmengen und Grundnotation]
\[
 A\subseteq B\ \Longleftrightarrow\
       \forall x\,(x\in A\Rightarrow x\in B),\qquad
 A\subsetneq B\ \Longleftrightarrow\ A\subseteq B\ \land\ A\ne B.
\]
$x\notin A$ verneint $x\in A$.
$\forall x\in A\,\varphi(x)$ und $\exists x\in A\,\varphi(x)$
sind beschränkte Quantoren. $\exists!x\,\varphi(x)$ bedeutet,
daß genau eine Menge $x$ die Formel erfüllt.
\end{definition}

\section{Die Axiome und Axiomenschemata}
\label{b03guide:sec:axioms}

Die folgende Liste legt die gesamte mengentheoretische Grundlage fest.
Die Zeichen $\varnothing$, $\{a,b\}$, $\bigcup A$ und
$\mathcal P(A)$ bezeichnen jeweils die durch das betreffende Axiom
und Extensionalität eindeutig bestimmte Menge.
Die Einzelmenge $\{a\}$ ist $\{a,a\}$; $A\cup B$ ist
$\bigcup\{A,B\}$. Diese Abkürzungen werden im folgenden verwendet.
\cite{pila,koepke}

\begin{axiom}[Extensionalität]
\[
 \forall A\,\forall B\,
 \bigl(\forall x\,(x\in A\Longleftrightarrow x\in B)
          \Rightarrow A=B\bigr).
\]
\end{axiom}

\begin{axiom}[Leere Menge]
\[
 \exists O\,\forall x\,(x\notin O).
\]
Diese Menge wird mit $\varnothing$ bezeichnet.
\end{axiom}

\begin{axiom}[Paarmenge]
\[
 \forall a\,\forall b\,\exists C\,\forall x\,
       (x\in C\Longleftrightarrow x=a\ \lor\ x=b).
\]
Diese Menge ist die ungeordnete Paarmenge $\{a,b\}$.
\end{axiom}

\begin{axiom}[Vereinigung]
\[
 \forall A\,\exists U\,\forall x\,
       (x\in U\Longleftrightarrow \exists Y\in A\,(x\in Y)).
\]
Diese Menge wird mit $\bigcup A$ bezeichnet.
\end{axiom}

\begin{axiom}[Potenzmenge]
\[
 \forall A\,\exists P\,\forall X\,
       (X\in P\Longleftrightarrow X\subseteq A).
\]
Diese Menge wird mit $\mathcal P(A)$ bezeichnet.
\end{axiom}

\begin{axiom}[Aussonderungsschema]
Für jede Formel $\varphi(x)$ mit festen Parametern gilt
\[
 \forall A\,\exists B\,\forall x\,
 (x\in B\Longleftrightarrow x\in A\ \land\ \varphi(x)).
\]
Die ausgesonderte Menge wird mit $\{x\in A\mid\varphi(x)\}$
bezeichnet.
\end{axiom}

\begin{axiom}[Ersetzungsschema]
Für jede Formel $\varphi(x,y)$ mit festen Parametern gilt
\[
 \forall A\,\Bigl[
 \bigl(\forall x\in A\,\exists!y\,\varphi(x,y)\bigr)
 \Rightarrow
 \exists B\,\forall y\,
 \bigl(y\in B\Longleftrightarrow
                 \exists x\in A\,\varphi(x,y)\bigr)\Bigr].
\]
Die Eindeutigkeit ist eine Voraussetzung an die Formel.
Ein Funktionsbegriff wird für dieses Schema nicht eingeführt.
\end{axiom}

\begin{axiom}[Unendlichkeit]
\[
 \exists I\,\bigl(\varnothing\in I\ \land\
             \forall x\in I\,(x\cup\{x\}\in I)\bigr).
\]
Hier wird nur das Axiom festgehalten; die natürlichen Zahlen werden
in Band~10 entwickelt.
\end{axiom}

\begin{axiom}[Fundierung oder Regularität]
\[
 \forall A\,\Bigl(A\ne\varnothing\Rightarrow
  \exists a\in A\,\neg\exists z\,(z\in a\ \land\ z\in A)\Bigr).
\]
\end{axiom}

Die leere Menge und das Aussonderungsschema können in anderen
Axiomatisierungen teilweise aus den übrigen Annahmen hergeleitet
werden. Ihre ausdrückliche Aufnahme entspricht hier der Darstellung
in Band~3. Zur Liste gehört kein Auswahlaxiom.
